\section{The Topo-Kinematic Circuit Identity}
\label{sec:topo_kinematic}

The Vacuum Circuit Analysis (VCA) framework rests upon a single, exact dimensional isomorphism between continuum spatial mechanics and electrical network theory. This section derives the mapping rigorously, verifies it through independent thermodynamic consistency checks, and tabulates the complete translation dictionary.

\subsection{Defining the Topological Conversion Constant}

The fundamental bridge between electrical and mechanical domains is the \textbf{Topological Conversion Constant} $\xi_{topo}$, defined as the ratio of the elementary charge to the spatial lattice pitch:
\begin{resultbox}{Topological Conversion Constant}
\begin{equation}
    \xi_{topo} \equiv \frac{e}{\ell_{node}} = \frac{1.602 \times 10^{-19} \text{ C}}{3.862 \times 10^{-13} \text{ m}} \approx 4.149 \times 10^{-7} \text{ C/m}
    \label{eq:xi_topo_vca}
\end{equation}
\end{resultbox}

This constant carries dimensions of $[\text{C/m}]$ and encodes the physical fact that, in a discrete lattice with pitch $\ell_{node}$, one quantum of electrical charge corresponds to one quantum of spatial displacement. Neither $e$ nor $\ell_{node}$ is adjustable; both are fixed by the lattice axioms (Axioms 1--2).

\subsection{Deriving the Six-Row Translation Table}

Each row of the Topo-Kinematic identity is derived by substituting $\xi_{topo}$ into the standard SI definition of the electrical quantity and isolating the resulting mechanical observable.

\paragraph{Row 1: Charge $\leftrightarrow$ Displacement.}
Electrical charge is the time integral of current. Substituting the kinematic mapping for current ($I \equiv \xi_{topo}\, v$):
\begin{resultbox}{Charge--Displacement Identity}
\begin{equation}
    Q = \int I \, dt = \xi_{topo} \int v \, dt = \xi_{topo}\, x
    \label{eq:charge_displacement}
\end{equation}
\end{resultbox}
\noindent Dimensional check: $[\text{C/m}] \times [\text{m}] = [\text{C}]$. \checkmark

\paragraph{Row 2: Current $\leftrightarrow$ Velocity.}
Differentiating the charge identity with respect to time:
\begin{resultbox}{Current--Velocity Identity}
\begin{equation}
    I = \frac{dQ}{dt} = \xi_{topo}\, \frac{dx}{dt} = \xi_{topo}\, v
    \label{eq:current_velocity}
\end{equation}
\end{resultbox}
\noindent The hardware velocity limit $v_{max} = c$ maps to a maximum circuit current:
\begin{equation}
    I_{max} = \xi_{topo}\, c = 4.149 \times 10^{-7} \times 2.998 \times 10^{8} \approx 124.4 \text{ A}
\end{equation}
This is the absolute slew-rate limit of the discrete lattice. No topological current can exceed $124.4$ A because no spatial velocity can exceed $c$.

\paragraph{Row 3: Voltage $\leftrightarrow$ Force.}
Voltage is the energy per unit charge ($V = W/Q$). Substituting $Q = \xi_{topo}\, x$ and $W = \int F\, dx$:
\begin{resultbox}{Voltage--Force Identity}
\begin{equation}
    V = \frac{W}{Q} = \frac{\int F\, dx}{\xi_{topo}\, x} \quad \Longrightarrow \quad V = \xi_{topo}^{-1}\, F
    \label{eq:voltage_force}
\end{equation}
\end{resultbox}
\noindent The maximum lattice force is the electromagnetic string tension $T_{EM} = m_e c^2 / \ell_{node} \approx 0.212$ N. The corresponding voltage is:
\begin{equation}
    V_{snap} = \xi_{topo}^{-1}\, T_{EM} = \frac{0.212}{4.149 \times 10^{-7}} \approx 511{,}000 \text{ V} = 511 \text{ kV}
\end{equation}
This recovers $V_{snap} = m_e c^2 / e$ exactly, confirming the identity.

\paragraph{Row 4: Inductance $\leftrightarrow$ Mass.}
Inductance opposes changes in current ($V = L\, dI/dt$). Substituting the voltage and current identities:
\begin{resultbox}{Inductance--Mass Identity}
\begin{equation}
    V = L \frac{dI}{dt} \quad \Longrightarrow \quad \xi_{topo}^{-1} F = L\, \xi_{topo} \frac{dv}{dt} = L\, \xi_{topo}\, a
    \quad \Longrightarrow \quad L = \xi_{topo}^{-2}\, m
    \label{eq:inductance_mass}
\end{equation}
\end{resultbox}
\noindent For the electron: $L_e = \xi_{topo}^{-2}\, m_e = m_e / \xi_{topo}^{2} \approx 5.292 \times 10^{-18}$ H ($\approx 5.3$ aH). This is the equivalent lumped inductance of a single electron's rest mass.

\paragraph{Row 5: Capacitance $\leftrightarrow$ Compliance.}
Capacitance stores charge per unit voltage ($C = Q/V$). Substituting:
\begin{resultbox}{Capacitance--Compliance Identity}
\begin{equation}
    C = \frac{Q}{V} = \frac{\xi_{topo}\, x}{\xi_{topo}^{-1}\, F} = \xi_{topo}^{2} \frac{x}{F} = \xi_{topo}^{2}\, \kappa
    \label{eq:capacitance_compliance}
\end{equation}
\end{resultbox}
\noindent where $\kappa = x/F = 1/k$ is the mechanical compliance (inverse spring constant). Dielectric breakdown occurs when the lattice displacement exceeds its absolute yield limit ($x > \ell_{node}$), which is structurally isomorphic to capacitor voltage exceeding the breakdown threshold.

\paragraph{Row 6: Resistance $\leftrightarrow$ Viscosity.}
Resistance dissipates power as $P = I^2 R$. Substituting:
\begin{resultbox}{Resistance--Viscosity Identity}
\begin{equation}
    R = \frac{V}{I} = \frac{\xi_{topo}^{-1} F}{\xi_{topo}\, v} = \xi_{topo}^{-2} \frac{F}{v} = \xi_{topo}^{-2}\, \eta
    \label{eq:resistance_viscosity}
\end{equation}
\end{resultbox}
\noindent where $\eta = F/v$ is the mechanical drag coefficient (units: kg/s). When the dielectric saturates and the vacuum enters the zero-impedance slipstream ($\eta \to 0$), the equivalent circuit resistance drops to zero---the hallmark of a superconductor.

\subsection{Complete Translation Dictionary}

Table~\ref{tab:topo_kinematic} consolidates the six identities into a single reference dictionary.

\begin{table}[H]
\centering
\caption{The Topo-Kinematic Circuit Identity. Every row is verified dimensionally and numerically from $\xi_{topo} = e/\ell_{node} \approx 4.149 \times 10^{-7}$ C/m.}
\label{tab:topo_kinematic}
\small
\begin{tabular}{@{}l l c l c@{}}
\toprule
\textbf{Electrical} & \textbf{Mechanical} & \textbf{Mapping} & \textbf{Units Check} & \textbf{Eq.} \\
\midrule
Charge $Q$        & Displacement $x$ & $Q = \xi\, x$             & $[\text{C/m}][\text{m}] = [\text{C}]$          & \ref{eq:charge_displacement} \\
Current $I$       & Velocity $v$     & $I = \xi\, v$             & $[\text{C/m}][\text{m/s}] = [\text{A}]$        & \ref{eq:current_velocity} \\
Voltage $V$       & Force $F$        & $V = \xi^{-1} F$          & $[\text{m/C}][\text{N}] = [\text{V}]$          & \ref{eq:voltage_force} \\
Inductance $L$    & Mass $m$         & $L = \xi^{-2} m$          & $[\text{m}^2/\text{C}^2][\text{kg}] = [\text{H}]$ & \ref{eq:inductance_mass} \\
Capacitance $C$   & Compliance $\kappa$ & $C = \xi^{2} \kappa$   & $[\text{C}^2/\text{m}^2][\text{m/N}] = [\text{F}]$ & \ref{eq:capacitance_compliance} \\
Resistance $R$    & Viscosity $\eta$ & $R = \xi^{-2} \eta$       & $[\text{m}^2/\text{C}^2][\text{kg/s}] = [\Omega]$  & \ref{eq:resistance_viscosity} \\
\bottomrule
\end{tabular}
\end{table}

\subsection{Self-Consistency Verification}

\paragraph{Work--Energy Theorem.}
The physical work done to charge a capacitor is $W = \int V\, dQ$. Substituting the Topo-Kinematic identities:
\begin{equation}
    W = \int (\xi_{topo}^{-1}\, F)(\xi_{topo}\, dx) = \int F\, dx
\end{equation}
The $\xi_{topo}$ factors cancel identically, recovering the mechanical work integral with no residual scaling. This proves the mapping is dimensionally exact and energy-conserving.

\paragraph{Impedance Cross-Check.}
The characteristic impedance of free space maps as:
\begin{equation}
    Z_{mech} = \xi_{topo}^{2}\, Z_0 = (4.149 \times 10^{-7})^2 \times 376.73 \approx 6.485 \times 10^{-11} \text{ kg/s}
\end{equation}
This is independently verifiable as the acoustic impedance of a single lattice node: $Z_{acoustic} = \rho_{bulk}\, c\, \ell_{node}^2 \approx 7.91 \times 10^6 \times 3.0 \times 10^8 \times (3.86 \times 10^{-13})^2 \approx 3.5 \times 10^{-10}$ kg/s. The factor-of-$\sim 5$ difference arises from the porosity correction ($p_c = 8\pi\alpha \approx 0.183$), confirming that the scaling is structurally consistent within the geometric packing fraction of the lattice.